\documentclass[12pt]{article}

\input{../../_assets/preambulo.tex}


\begin{document}

    % 1. Foto de fondo
    % 2. Título
    % 3. Encabezado Izquierdo
    % 4. Color de fondo
    % 5. Coord x del titulo
    % 6. Coord y del titulo
    % 7. Fecha

    
    \input{../../_assets/portada}
    \portadaExamen{ffccA4.jpg}{Ecuaciones\\Diferenciales II\\Examen XXI}{Ecuaciones Diferenciales II. Examen XXI}{MidnightBlue}{-8}{28}{2026}{}

    \begin{description}
        \item[Asignatura] Ecuaciones Diferenciales II.
        \item[Curso Académico] 2022/2023.
        \item[Grado] Doble Grado en Ingeniería Informática y Matemáticas.
        \item[Grupo] Único.
        \item[Profesor] Rafael Ortega Ríos.
        \item[Descripción] Convocatoria Extraordinaria.
        \item[Fecha] 11 de julio de 2023.
        % \item[Duración] ---.
    \end{description}
    \newpage


    % ------------------------------------
    
    \begin{ejercicio}
        Para cada $n \geq 1$ la función $f_n : \bb{R} \to \bb{R}$ se define como la característica del intervalo $\left[\nicefrac{1}{n}, 1\right]$; es decir
        \[
        f_n(t) = \begin{cases}
            1 & \text{si } t \in \left[\nicefrac{1}{n}, 1\right] \\
            0 & \text{si } t \in \bb{R} \setminus \left[\nicefrac{1}{n}, 1\right]
        \end{cases}
        \]
        ¿Hay convergencia puntual de la sucesión $\{f_n\}_{n \geq 1}$? ¿Hay convergencia uniforme?
    \end{ejercicio}

    \begin{ejercicio}
        Se considera el problema de valores iniciales
        \[
        \dot{x} = f(x), \quad x(t_0) = x_0
        \]
        donde
        \[
        f : \bb{R} \to \bb{R}, \quad f(x) = \begin{cases}
            2x & \text{si } x \geq 0 \\
            3x & \text{si } x < 0
        \end{cases}
        \]
        ¿Hay unicidad? Encuentra la solución maximal.
    \end{ejercicio}

    \begin{ejercicio}
        Dados $t_0, x_0 \in \bb{R}$, encuentra el intervalo maximal de la solución del problema
        \[
        \dot{x} = x^5, \quad x(t_0) = x_0.
        \]
    \end{ejercicio}

    \begin{ejercicio}
        Construye una retracción de $\bb{R}^2$ en el conjunto
        \[
        A = \left\{(x,y) \in \bb{R}^2 : |x| + |y| \leq 3\right\}.
        \]
    \end{ejercicio}

    \begin{ejercicio}
        La solución del problema de Cauchy
        \[
        \ddot{\theta} + 9 \operatorname{sen} \theta = 0, \quad \theta(0) = \theta_0, \quad \dot{\theta}(0) = \omega_0
        \]
        se denota por $\theta(t; \theta_0, \omega_0)$. Calcula, si existe, $\deld{\theta}{\theta_0}(t; 3\pi, 0)$.
    \end{ejercicio}

    \begin{ejercicio}~
        \begin{enumerate}[label=\roman*)]
            \item Dado un parámetro $\veps > 0$, $x(t, \veps)$ es la solución del problema
            \[
            \dot{x} = \frac{\operatorname{sen}(\veps x)}{\veps} + e^t, \quad x(0) = \pi.
            \]
            Calcula, si existe, $\lim\limits_{\veps \to 0^+} x(2\pi, \veps)$.
            \item Dado un parámetro $\veps > 0$, $x(t, \veps)$ es la solución del problema
            \[
            \veps \dot{x} = e^x \operatorname{sen} t, \quad x(0) = \pi.
            \]
            Calcula, si existe, $\lim\limits_{\veps \to 0^+} x(2\pi, \veps)$.
        \end{enumerate}
    \end{ejercicio}

    \newpage
    \setcounter{ejercicio}{0} % Reiniciar contador de ejercicios
    \noindent
    % \textbf{Solución.}

\end{document}